\section{Introduction}%
\label{sec:introduction}
%Each comment introduces 1-2 sentences 
% Context 
Knowledge graphs (KGs) have seen increased usage alongside large language
models (LLMs), enhancing LLMs with external knowledge for tasks such as
retrieval-augmented generation (RAG)~\cite{Xu2024RAGKGLinkedin}. 
The KGs required for LLM applications
	\todo{KGs are not generally "required for LLM applications" (but for different types of RAG applications which, in turn, also use LLMs). \\ In addition to this minor issue, a more general comment is that this paragraph puts too much emphasis on LLMs. I would tone it done and start a bit broader regarding the use cases of KGs, and then you can mention KGs for RAG (and LLMs) as a more recent group of use cases. At the moment, the paragraph is almost more about LLMs and KGs for LLMs and not so much about KGs per se.}
may already exist or may need to be
generated from existing (semi-)structured
	\todo{Besides the issue with the closing parenthesis, what are "(semi-)structured data sources"? Do you mean "structured or semi-structured data sources"?}
data sources such as 
tabular data in CSV formats
	\todo{It is not clear why you are using plural here ("CSV format\textbf{s}").}
or hierarchical data in JSON.
	\todo{"hierarchical data " sounds weird (how can data be hierarchical?). etter say: "nested data objects"}
Naturally, empowering LLMs with the vast amount of existing (semi-)structured
data enables the reuse of existing data-generating infrastructures, 
and eliminating the additional costs of setting up new KG-generating sources.
One approach to generate KG
	\todo{As a general comment, please pay attention to correct usage (or omission) of articles, and plural versus singular versions of nouns. Here it needs to be either "to generate a KG" or "to generate KGs"}
from existing (semi-)structured data sources,
	\todo{Why do you have a comma here? I have noticed also earlier that you are putting commas in strange places. Please try to avoid that. Here is a prompt for ChatGPT based on which you can check your commas: \\
	Consider the following sentence, which I will give in quotation marks, is it correct to have the comma in this sentence?
	"One approach to generate KG from existing (semi-)structured data sources, is through the usage of declarative mapping languages."}
is
through the usage of declarative mapping
languages~\cite{VanAssche2022DeclarativeRDFgraph}.

Declarative mapping languages specify the process
	\todo{No, such languages do not specify any process. They allow users to specify/define mappings.}
for generating KGs from
(semi-)structured data based on a mapping document written in the respective
language. 
RDF Mapping Language
	\todo{\textbf{The} RDF Mapping Language ...}
(RML)~\cite{dimou_ldow_2014} is one example of
such a declarative mapping language.
Usage of RML, to generate KGs in RDF format from heterogeneous data sources,
	\todo{While this sentence, with its commas is not incorrect, it sounds unnatural (a native English speaker would never write it this way). Remove the commas and replace "to generate" by "for generating"}
has gained traction in recent years across diverse domains such as COVID
research~\cite{Sakor2023KG4Covid},
	%European
railway networks~\cite{rojas2021leveraging}, and
incident management in ICT systems~\cite{Tailhardat2024NORIA}. 
Numerous
	\todo{"numerous" means many. Of course, it is subjective what many are, but in this context I would not use this word. Better: "A number of .." or "Several .."}
RML-based mapping engines have been developed to generate KG
	\todo{KGs}
from 
heterogeneous data, each with its unique approach to execute the mapping process
(i.e. generating KG from existing semi-structured data~\cite{VanAssche2022DeclarativeRDFgraph}). 
For example, FlexRML~\cite{Freund2024FlexRML} focuses on the generation of
	\todo{KGs}
KG for 
resource-constrained IoT devices, while SDM-RDFizer~\cite{Iglesias2022ScalingKnowledgeGraph} 
aims for an efficient and scalable KG generation through novel 
optimization techniques and data structures. 

% what’s the problem
Current state-of-the-art mapping engines exhibit varying operational semantics
when executing the mapping processes specified within an RML document,
	\todo{Again, an RML document does not specify any process.}
which is the consequence of the lack of formalizations in RML. 
The lack of formalizations not only complicates the verification
of the correctness and the study of the expressiveness of the language, it also
makes it challenging to optimize the mapping process described by the RML
document.
Furthermore, the formalizations should be mapping language-agnostic,
	\todo{By writing it this way around, this sentence here immediately triggers the question: Why? (Why should the formalizations for RML be language agnostic?)}
and 
provide operational semantics
	\todo{I don't think you really mean operational semantics here.}
of a general mapping process. 
	\todo{What is "a general mapping process"?}
%To the best of our knowledge, there is only one study that provided 
%formalizations to RML by E. Iglesias et.al.~\cite{Iglesias2022ScalingKnowledgeGraph}.
%However, their formalizations suffer from being dependent upon the syntaxes 
%of RML and limited to the context of optimizing RML document.
%Other studies also provide formalizations of declarative mapping languages such
%as XSPARQL~\cite{Lopes2011XSPARQL}, and Façade-X~\cite{asprino2023knowledge}
%but remain limited to the language that they use. 
%There are currently no studies on language-agnostic formalizations of 
%mapping processes. 
% why it’s an interesting problem
A language-agnostic formalization of the mapping process paves the way for a
unified mapping engine capable of working with the user's choice of language,
ensuring verified correctness and supporting optimizations across different
mapping languages.

\todo[inline]{The previous two paragraphs are too RML centric! And they do not clear enough point out the general gap. \\ The message of the first of these two paragraphs should be that there exists several declarative mapping approaches to define and to create KGs (including RML) and that these approaches are being used for many applications. Or perhaps even better the other way around: 
Declarative mapping approaches are being used for many applications and there are several such approaches (including RML) and corresponding systems/engines. The message of the second paragraph should then be that all these approaches lack a well-defined formal foundation and that this is a problem because ... . (And the part about the advantages of a language-agnostic formalism should not be here because this is not related to the problem, as the problem is the lack of any formal foundation, language specific or language agnostic. So, move the part about these advantages to the next paragraph.)}

% what’s my solution
In this work, we tackled
	\todo{The paper should we written in present tense.}
the challenge
	\todo{It is not clear why that is a challenge. After changing the previous paragraph as suggested in my previous comment, you can better start here by saying that, in this paper, we fill this gap.}
of providing mapping language-agnostic
formalizations in two steps.
First, we provide formal definitions and propose a set of algebraic mapping
operators which can be composed to formally describe the mapping processes. 
In the second and last step, we provide an algorithm for translating RML documents
	\todo{Introducing our work in terms of these two steps is, again, too RML centric. In fact, the first three sentences here can be read as: The authors introduce a language-agnostic formalization to define RML. Or perhaps even: The authors introduce a language-agnostic formalization of RML. Yet, RML is not the main thing here! We are making multiple contributions, the primary one being the introduction of formal approach to define mappings (which is language agnostic and you can emphasize the advantages of this fact here) and another one being that we show that, with our approach, it is possible to capture every mapping that can be defined using RML. As a side effect of this second contribution, we are actually providing a formal semantics for RML, which can be seen as another contribution. (And then we also have algebraic equivalences as another contribution.)}
to an expression comprising these algebraic mapping operators,
thereby formalizing RML's operational semantics and demonstrating that our 
language-agnostic algebra can serve as the foundation for formalizing
various mapping languages.
	\todo{This claim about "various mapping languages" is too strong, considering what we do in this paper.}
% what comes after my solution
The set of algebraic mapping operators leads to a unified model for 
the mapping process of generating KG from (semi-)structured data, 
optimization techniques applicable to all implementations of our 
algebraic operators, and an algebraic mapping engine that could utilize 
multiple declarative mapping languages without any changes to the 
implementation.


% paper outline
The paper is outlined as follows.
	 \todo{I personally find it quite useless and a waste of valuable space to have such a paragraph that simply lists all sections with a brief mention of their respective content. If you think about it, do you ever read these paragraphs?? What I typically do instead, is to integrate the section references into the text that introduces the contributions. As concrete examples, you may look at the following three papers:
	 %https://olafhartig.de/files/ChengHartig_FedQPL_iiWAS2020_Extended.pdf
	 %https://olafhartig.de/files/ChengHartig_OptPlus_2019.pdf
	 %https://olafhartig.de/files/ChengHartig_LinGBM_WISE2022_AcceptedManuscript.pdf
	 }
Section~\ref{sec:related_works} provides an overview of related work in the domain of  
formalizing declarative mapping languages. 
Section~\ref{sec:data_model} introduces the notations and symbols we 
used to provide the formal definitions of algebraic mapping operators. 
Section~\ref{sec:mapping_algebra} contains the formal definitions of 
our algebraic mapping operators and their semantics. 
In Section~\ref{sec:rml_translation} we provide our translation algorithm
for transforming RML document into an expression composed of our algebraic 
mapping operators, thereby formalizing RML. 
Section~\ref{sec:equivalences} introduces a few equivalences and proofs for 
optimizations of the generated mapping plan. 
Finally, we conclude this paper and provides possible future works utilizing
our algebraic mapping operators in Section~\ref{sec:conclusion}. 

